%! Author = nikitaivanov
%! Date = 12/04/2026

\documentclass[a4paper,14pt]{extarticle}

\usepackage[russian]{babel}
\usepackage{caption}
\usepackage{enumitem}
\usepackage{fontspec}
\usepackage[left=3cm, right=1.5cm, top=2cm, bottom=2cm]{geometry}
\usepackage{graphicx}
\usepackage{mfirstuc}
\usepackage{pdflscape}
\usepackage{setspace}
\usepackage{titlesec}
\usepackage{longtable}
\usepackage{fancybox}
\usepackage{tocloft}
\usepackage{titletoc}
\usepackage{listings}
\usepackage{xcolor}
\usepackage{multirow}
\usepackage{minted}
\usepackage[hyphens]{url}
\usepackage{indentfirst}
\usepackage{ragged2e}
\usepackage[hidelinks]{hyperref}

\Urlmuskip=0mu plus 1mu
\setlength{\emergencystretch}{3em}
\makeatletter
\providecommand{\babel@toc}[2]{}
\makeatother
\pdfstringdefDisableCommands{%
  \def\uppercase#1{#1}%
  \def\makefirstuc#1{#1}%
  \def\normalsize{}%
  \def\textnormal#1{#1}%
  \def\textbf#1{#1}%
}

\definecolor{codegreen}{rgb}{0,0.6,0}
\definecolor{codegray}{rgb}{0.5,0.5,0.5}
\definecolor{codepurple}{rgb}{0.58,0,0.82}
\definecolor{backcolour}{rgb}{0.95,0.95,0.92}

\lstdefinestyle{mystyle}{
  backgroundcolor=\color{backcolour},
  commentstyle=\color{codegreen},
  keywordstyle=\color{magenta},
  numberstyle=\tiny\color{codegray},
  stringstyle=\color{codepurple},
  basicstyle=\ttfamily\footnotesize,
  breakatwhitespace=false,
  breaklines=true,
  captionpos=b,
  keepspaces=true,
  numbers=left,
  numbersep=5pt,
  showspaces=false,
  showstringspaces=false,
  showtabs=false,
  tabsize=2
}

\lstdefinelanguage{JavaScript}{
  keywords={break, case, catch, const, continue, debugger, default, delete, do, else, finally, for, function, if, in, instanceof, let, new, return, switch, this, throw, try, typeof, var, void, while},
  sensitive=true,
  comment=[l]{//},
  morecomment=[s]{/*}{*/},
  morestring=[b]",
  morestring=[b]'
}

\lstset{style=mystyle}

\setmainfont{Times New Roman}
\setmonofont{Courier New}
\renewcommand{\baselinestretch}{1.5}
\setlength{\parindent}{1.25cm}

\addto\captionsrussian{\renewcommand{\figurename}{Рисунок}}
\renewcommand{\tablename}{Таблица}
\DeclareCaptionLabelSeparator{dash}{ -- }
\captionsetup{labelsep=dash}
\captionsetup[figure]{name=Рисунок, labelsep=dash}

\renewcommand{\thesection}{\textnormal\textbf{\normalsize\arabic{section}}}
\renewcommand{\thesubsection}{\thesection.\normalsize\arabic{subsection}}
\renewcommand{\thefigure}{\arabic{figure}}
\renewcommand{\thetable}{\arabic{table}}

\titleformat{\section}
    {\centering\normalsize\bfseries}
    {ГЛАВА \thesection.}
    {0.1em}
    {\uppercase}

\titleformat{\subsection}
    {\normalsize\bfseries}
    {\thesubsection}
    {0.5em}
    {\makefirstuc}

\renewcommand\cfttoctitlefont{\hfill\normalfont\bfseries}
\renewcommand\cftaftertoctitle{\hfill\mbox{}}
\renewcommand\contentsname{СОДЕРЖАНИЕ}

\setcounter{tocdepth}{2}
\setcounter{secnumdepth}{2}

\newcommand{\withoutnumber}[1]{%
    \clearpage
    \section*{\centering\normalsize\uppercase{#1}}
    \addcontentsline{toc}{section}{\normalsize \uppercase{#1}}
}

\newcommand{\withnumber}[1]{%
    \clearpage
    \section{\uppercase{#1}}
}

\newcommand{\mysubsection}[1]{%
    \subsection{\normalsize\makefirstuc{#1}}
}

\newcommand{\illustration}[3]{%
    \begin{figure}[htbp]
        \centering
        \includegraphics[width=#2\textwidth]{#1}
        \captionsetup{font=singlespacing}
        \caption{\makefirstuc{#3}}
        \label{fig:#1}
    \end{figure}
}

\newcommand{\image}[2]{%
    \begin{figure}[htbp]
        \centering
        \includegraphics[width=\textwidth]{#1}
        \captionsetup{font=singlespacing}
        \caption{\capitalisewords{#2}}
        \label{fig:#1}
    \end{figure}
}

\newcommand{\unlabeledimage}[3]{%
    \begin{figure}[htbp]
        \centering
        \includegraphics[width=#2\textwidth]{#1}
        \captionsetup{font=singlespacing}
        \caption{\makefirstuc{#3}}
    \end{figure}
}

\newcommand{\placeholderimage}[2]{%
    \begin{figure}[htbp]
        \centering
        \fbox{%
            \parbox[c][6cm][c]{0.85\textwidth}{%
                \centering
                Шаблон рисунка\\[0.5em]
                Файл: \texttt{#1}
            }%
        }
        \captionsetup{font=singlespacing}
        \caption{\makefirstuc{#2}}
        \label{fig:placeholder:#1}
    \end{figure}
}

\newcommand{\tablecaption}[2]{%
    \captionsetup{
        justification=raggedright,
        singlelinecheck=false,
        font=singlespacing,
        labelsep=dash
    }
    \caption{\makefirstuc{#1}}
    \label{tab:#2}
}

\newenvironment{reporttable}[1]
{%
    \begin{table}[htbp]
    \centering
    \footnotesize
    \captionsetup{font=singlespacing}
    #1
}
{%
    \end{table}
}

\newenvironment{reportlongtable}[1]
{%
    \begin{footnotesize}
    \begin{longtable}{#1}
}
{%
    \end{longtable}
    \end{footnotesize}
}

\newcommand{\codeexample}[3]{%
    \lstinputlisting[language=#2,caption={#3}]{#1}
}

\newcommand{\appendixsection}[1]{%
    \withoutnumber{#1}
}

\begin{document}

\begin{titlepage}
\thispagestyle{empty}
\begin{center}
    \textbf{Министерство науки и высшего образования Российской Федерации}\\
    \textbf{Федеральное государственное автономное образовательное учреждение высшего образования}\\
    \textbf{«Национальный исследовательский университет ИТМО»}\\
    (Университет ИТМО)

    \vspace{1cm}

    \textbf{Факультет программной инженерии и компьютерной техники}\\
    \textbf{Образовательная программа «Веб-технологии»}

    \vspace{2cm}

    \textbf{\large О Т Ч Ё Т}\\
    \textbf{по дисциплине «Базы данных и знаний»}

    \vspace{1.5cm}

    \begin{flushleft}
    \textbf{Тема:} Проектирование и анализ схемы хранения композитных объектов, представляемых на основе pure operation-based CRDT

    \vspace{0.5cm}

    \textbf{Обучающийся:} Иванов Никита Денисович

    \vspace{0.5cm}

    \textbf{Группа:} P4110

    \vspace{0.5cm}

    \textbf{Срок выполнения:} апрель 2026
    \end{flushleft}

    \vfill

    Санкт-Петербург\\
    2026
\end{center}
\end{titlepage}

\pagenumbering{arabic}
\setcounter{page}{2}

\tableofcontents
\clearpage

\withoutnumber{Введение}

В рамках индивидуального задания рассматривается задача проектирования хранилища
для композитных объектов, состояние которых определяется журналом операций.
Исходная предметная область связана с прототипом коллаборативного календарного
приложения, поддерживающего оффлайн-режим, локальные изменения и последующую
синхронизацию между репликами.

Весь исходный код, конфигурационные файлы, материалы по экспериментам и
исходный текст отчёта хранятся в открытом репозитории GitHub:
\url{https://github.com/Nikita2672/databases_individual_task}.

Актуальность темы обусловлена тем, что при построении систем на основе CRDT
необходимо не только определить прикладную модель конфликтов и синхронизации,
но и подобрать такую структуру хранения, которая будет одновременно удобной
для записи операций, восстановления истории объекта и анализа производительности
типовых запросов.

Целью работы является сравнение трёх вариантов хранения одной и той же модели
операций: в реляционной СУБД PostgreSQL, в документной СУБД MongoDB и в
распределённой платформе YTsaurus.

Для достижения поставленной цели решаются следующие задачи:
\begin{itemize}[leftmargin=1.25cm]
    \item определить логическую структуру данных, описывающую реплики, объекты и операции;
    \item спроектировать схему хранения в PostgreSQL;
    \item спроектировать схему хранения в MongoDB;
    \item рассмотреть возможную модель хранения в YTsaurus;
    \item проанализировать сильные и слабые стороны каждого подхода;
    \item подготовить основу для дальнейшего анализа индексов, планов выполнения запросов и ограничений каждой платформы.
\end{itemize}

\withnumber{Постановка задачи и модель данных}

Репозиторий выделен как отдельная инфраструктурная часть индивидуального задания
и предназначен для исследования нескольких вариантов хранения одного и того же
контракта операций CRDT. Практический интерес работы состоит не только в выборе
СУБД, но и в сравнении разных способов организации одного и того же operation log.

\mysubsection{Состав репозитория}

В текущем состоянии репозиторий включает следующие основные части:
\begin{itemize}[leftmargin=1.25cm]
    \item каталог \texttt{postgresql} с \texttt{Docker Compose}, миграцией схемы и материалами по анализу запросов;
    \item каталог \texttt{mongodb} с \texttt{Docker Compose}, инициализацией коллекции и материалами по \texttt{executionStats};
    \item каталог \texttt{YTsaurus} с Java-классами генерации и сценариями запросов для распределённого варианта;
    \item каталог \texttt{report} с исходным текстом данного отчёта.
\end{itemize}

Таким образом, работа организована как воспроизводимое исследовательское
окружение: для PostgreSQL и MongoDB используются локально поднимаемые
изолированные инсталляции, а для YTsaurus рассматривается отдельный
распределённый контур выполнения.

Следует отдельно отметить, что вариант на \texttt{YTsaurus} не развёртывался
локально в пределах данного репозитория. Для него было принято решение
использовать уже существующий кластер в рабочей инфраструктуре VK, на котором
и проводились все тесты. Поэтому размещённые в каталоге \texttt{YTsaurus}
Java-файлы следует рассматривать прежде всего как листинги фактически
использованного кода и иллюстрацию реализованных сценариев, а не как
самодостаточный локально воспроизводимый модуль, который можно запустить в
отрыве от соответствующего инфраструктурного окружения.

\mysubsection{Фактическая структура операции}

Согласно описанию в \texttt{README.md} и специализированных подкаталогах,
центральной единицей прикладной модели является операция со следующими полями:
\begin{itemize}[leftmargin=1.25cm]
    \item \texttt{opId};
    \item \texttt{txId};
    \item \texttt{objectId};
    \item \texttt{replicaId};
    \item \texttt{clock};
    \item \texttt{action}.
\end{itemize}

Поле \texttt{clock} хранит причинно-следственную информацию.
Поле \texttt{action} содержит полезную нагрузку операции.
Во всех рассматриваемых вариантах используется одна и та же прикладная структура;
различается только физическая организация данных.

\withnumber{Реализация варианта на PostgreSQL}

В каталоге \texttt{postgresql} реализовано автономное окружение для
реляционного моделирования журнала CRDT-операций. Запуск выполняется через
\path{docker compose -f postgresql/compose.yaml up -d}, а физическая схема
создаётся миграцией \path{postgresql/initdb/001_schema.sql}.

\mysubsection{Назначение PostgreSQL-варианта}

PostgreSQL используется как базовая реляционная СУБД для хранения append-only
operation log. Такой вариант удобен для анализа, поскольку позволяет явно задать
структуру данных, внешние ключи, ограничения целостности и индексную стратегию,
а затем исследовать поведение запросов через \texttt{EXPLAIN ANALYZE}.

\mysubsection{Логическая и физическая схема}

Текущая реляционная схема включает три таблицы:
\begin{itemize}[leftmargin=1.25cm]
    \item \texttt{replicas} --- справочник логических реплик;
    \item \texttt{objects} --- справочник корневых CRDT-объектов;
    \item \texttt{operations} --- неизменяемый журнал операций.
\end{itemize}

Связи между сущностями заданы явно: одна реплика может породить много операций,
один объект может иметь много операций, каждая операция принадлежит ровно одной
реплике и ровно одному объекту. При этом отдельная таблица зависимостей между
операциями не вводится, поскольку причинность задаётся содержимым поля
\texttt{clock}.

Физическая схема повторяет логическую модель почти напрямую. Идентификаторы
хранятся как \texttt{TEXT}, что соответствует строковому формату библиотеки.
Поля \texttt{clock} и \texttt{action} сохраняются в \texttt{JSONB}, поскольку
они имеют полуструктурированную природу и на текущем этапе исследования не
декомпозируются в набор отдельных реляционных таблиц.

\illustration{../postgresql/resources/diagram.png}{0.8}{Логическая ER-диаграмма варианта хранения в PostgreSQL}

\illustration{../postgresql/resources/objects.png}{1.0}{Физическая схема таблиц PostgreSQL}

\mysubsection{Целостность и нормализация}

Схему PostgreSQL корректно считать частично нормализованной. Каркас из таблиц
\texttt{replicas}, \texttt{objects} и \texttt{operations} соответствует обычной
реляционной модели и обеспечивает:
\begin{itemize}[leftmargin=1.25cm]
    \item первичные ключи для всех основных сущностей;
    \item внешние ключи из \texttt{operations} в \texttt{replicas} и \texttt{objects};
    \item \texttt{CHECK}-ограничения, контролирующие, что \texttt{clock} и \texttt{action} являются JSON-объектами;
    \item проверку наличия поля \texttt{type} внутри \texttt{action}.
\end{itemize}

При этом доменная корректность содержимого \texttt{clock} и \texttt{action}
остаётся обязанностью прикладного уровня. Это означает, что PostgreSQL отвечает
за структурную и ссылочную целостность, а библиотека --- за корректную семантику
CRDT-операций.

\mysubsection{Основные сценарии запросов}

Из \texttt{postgresql/README.md} видно, что в реляционном варианте в фокусе
находятся выборки по операциям как самостоятельным записям. Базовые сценарии:
\begin{itemize}[leftmargin=1.25cm]
    \item получение полной истории операций для заданного \texttt{objectId};
    \item получение операций конкретной транзакции по \texttt{txId};
    \item получение операций, выпущенных конкретной репликой;
    \item чтение сохранённых векторных часов для конкретной операции;
    \item фильтрация по типу действия внутри \texttt{action}.
\end{itemize}

Этот набор запросов показывает, что PostgreSQL ориентирован не только на чтение
истории одного объекта, но и на глобальный анализ журнала операций.

\mysubsection{Фрагмент SQL-схемы}

Ниже приведён фактический фрагмент миграции, создающей основные таблицы и индексы
реляционного варианта:

\codeexample{../postgresql/initdb/001_schema.sql}{SQL}{Миграция схемы PostgreSQL}

\mysubsection{Преимущества и ограничения}

К основным преимуществам PostgreSQL-варианта относятся строгая схема хранения,
прозрачная ссылочная целостность, удобство аналитических выборок по
\texttt{opId}, \texttt{txId} и \texttt{objectId}, а также естественная поддержка
append-only журнала как множества независимых строк.

Ограничения реляционного варианта также явно видны. Во-первых, часть структуры
операции остаётся внутри \texttt{JSONB}, поэтому схема является лишь частично
нормализованной. Во-вторых, длинная история одного объекта требует сортировки
и может становиться заметной по стоимости. В-третьих, низкоселективные широкие
выборки могут приводить к менее эффективным планам, включая \texttt{Seq Scan}.

\withnumber{Реализация варианта на MongoDB}

В каталоге \texttt{mongodb} размещено автономное окружение документной СУБД
MongoDB. Запуск выполняется через \path{docker compose -f mongodb/compose.yaml up -d},
а структура коллекции и индексов задаётся скриптом
\path{mongodb/initdb/001_init.js}.

\mysubsection{Назначение MongoDB-варианта}

MongoDB рассматривается как объектно-центричная альтернатива PostgreSQL. Если
в реляционном варианте единицей хранения является отдельная операция, то здесь
основной единицей становится документ объекта со встроенным массивом операций.

\mysubsection{Выбранная документная модель}

Для текущего этапа принята модель ``один объект --- один документ со встроенной
историей операций''. В текущем варианте используется одна коллекция:
\begin{itemize}[leftmargin=1.25cm]
    \item \texttt{objects}.
\end{itemize}

Документ верхнего уровня содержит:
\begin{itemize}[leftmargin=1.25cm]
    \item \texttt{objectId};
    \item массив \texttt{operations}.
\end{itemize}

Каждый элемент массива \texttt{operations} соответствует библиотечному типу
\texttt{Operation} и содержит те же поля, что и в реляционном варианте. Таким
образом, меняется не состав данных, а единица агрегирования на уровне хранилища.

\illustration{../mongodb/resources/diagram.png}{0.95}{Структура документа объекта в MongoDB}

\mysubsection{Индексация и сценарии чтения}

В \texttt{mongodb/README.md} зафиксированы следующие направления индексации:
\begin{itemize}[leftmargin=1.25cm]
    \item уникальный индекс по \texttt{objectId};
    \item индекс по \texttt{operations.opId};
    \item индекс по \texttt{operations.replicaId};
    \item индекс по \texttt{operations.action.type}.
\end{itemize}

Это хорошо отражает специфику документной модели: индексируются вложенные поля
массива операций, а не top-level коллекция операций. Основные сценарии запросов
при такой организации следующие:
\begin{itemize}[leftmargin=1.25cm]
    \item поиск документа объекта по \texttt{objectId};
    \item чтение полной встроенной истории конкретного объекта;
    \item поиск объекта, содержащего операцию с заданным \texttt{opId};
    \item поиск объектов, в истории которых есть операции конкретной реплики;
    \item поиск по значениям \texttt{operations.action.type}.
\end{itemize}

\mysubsection{Преимущества и ограничения}

Сильные стороны MongoDB-варианта заключаются в естественном хранении
полуструктурированных полей \texttt{clock} и \texttt{action}, а также в том,
что полная история конкретного объекта может быть прочитана из одного документа
без отдельного журнала top-level операций.

Ограничения модели также существенны. Размер документа MongoDB ограничен
\texttt{16 MB}, очень длинная история объекта ухудшает устойчивость embedded-модели,
а глобальные выборки по операциям через весь набор документов менее естественны,
чем в PostgreSQL. В \texttt{README} также отмечено, что при дальнейшем росте
истории может потребоваться сегментация документа.

\mysubsection{Фрагмент инициализации MongoDB}

Ниже приведён фактический фрагмент скрипта, создающего коллекцию \texttt{objects}
и основные индексы документного варианта:

\codeexample{../mongodb/initdb/001_init.js}{JavaScript}{Скрипт инициализации MongoDB}

\withnumber{Вариант на YTsaurus}

\mysubsection{Общая характеристика платформы}

YTsaurus представляет собой отечественную open-source платформу для
распределённого хранения и обработки больших объёмов данных. В отличие от
классических СУБД, ориентированных преимущественно либо на транзакционные,
либо на аналитические сценарии, YTsaurus изначально проектируется как
многофункциональная инфраструктурная платформа, объединяющая распределённое
хранилище данных, средства пакетной обработки, механизмы выполнения
SQL-подобных запросов и поддержку динамических таблиц.

Архитектурно YTsaurus сочетает несколько уровней. Базовый уровень образует
распределённое хранилище и пространство метаданных \texttt{Cypress}, в котором
данные организуются в виде древовидного пространства имён. Поверх него
реализуются механизмы хранения статических и динамических таблиц, а также
вычислительная подсистема, поддерживающая \texttt{MapReduce}-обработку и более
высокоуровневые инструменты, такие как \texttt{YQL}, \texttt{CHYT} и
\texttt{SPYT}. Таким образом, YTsaurus следует рассматривать не как отдельную
узкоспециализированную базу данных, а как целостную экосистему для хранения,
трансформации и анализа данных в распределённой среде.

Практическое назначение YTsaurus связано с задачами, в которых требуется
одновременно обеспечить горизонтальную масштабируемость, отказоустойчивость,
изолированное сосуществование нескольких нагрузок и поддержку различных
способов доступа к данным. Платформа ориентирована на сценарии обработки
больших массивов событий, логов, транзакционных записей и иных наборов данных,
для которых важны не только объём хранения, но и возможность последующего
параллельного анализа.

С точки зрения настоящей работы интерес к YTsaurus обусловлен тем, что данная
платформа предлагает ещё один принципиально отличный контекст хранения данных
по сравнению с PostgreSQL и MongoDB. Если PostgreSQL в данном исследовании
рассматривается как реляционное хранилище глобального журнала операций, а
MongoDB --- как документное объектно-центричное хранилище встроенной истории,
то YTsaurus представляет собой распределённую платформу, в рамках которой
необходимо отдельно определить, какая модель представления operation-based
CRDT будет наиболее естественной: журнал операций, табличное представление,
ключ-значение или их комбинация.

\mysubsection{Экспериментальное окружение}

Практическая часть для YTsaurus выполнялась не в локальном контейнерном
окружении, а на реальных кластерах VK, используемых в инфраструктуре
рекомендательной платформы.

В отличие от PostgreSQL и MongoDB, где в работе можно было
опираться на \texttt{EXPLAIN ANALYZE} и \texttt{explain("executionStats")},
в рассматриваемом окружении YTsaurus не предоставлял API для аналогичного
плана выполнения запросов. Поэтому для YTsaurus удалось замерить только
внешнее время исполнения запросов и размер результирующих выборок, без
детального разбора внутреннего плана, выбора стратегий чтения и стоимости
отдельных стадий выполнения.

\mysubsection{Реализованная модель хранения}

В каталоге \texttt{YTsaurus/java\_files} размещены Java-классы, использовавшиеся
для генерации данных и выполнения простых сценариев чтения. Центральным
классом доступа к данным является \texttt{OperationDao}, работающий с одной
dynamic table по пути \path{//home/dev/n.d.ivanov/operations}. Каждая строка
таблицы соответствует одной неизменяемой операции \texttt{op-based CRDT}.

Содержательно эта модель близка к PostgreSQL-варианту: единицей хранения
остается отдельная операция, а не агрегированный документ объекта. В записи
сохраняются:
\begin{itemize}[leftmargin=1.25cm]
    \item составной ключ \texttt{objectId}, \texttt{replicaId}, \texttt{txId}, \texttt{opId};
    \item поле \texttt{clock} как снимок векторных часов;
    \item поле \texttt{action} как сериализованное CRDT-действие.
\end{itemize}

Порядок полей в составном ключе выбран не случайно. Он оптимизирован под
типовой доступ по \texttt{objectId}, затем по \texttt{replicaId} и
\texttt{txId}. Тем самым YTsaurus-вариант в данной реализации также моделирует
глобальный журнал операций, но уже поверх распределённой dynamic table.

\mysubsection{Генератор данных и сценарии запросов}

Класс \texttt{OperationGenerator} формирует синтетический набор операций в
домене календарных событий. Генерация выполняется для нескольких объектов и
реплик, при этом для каждой операции создаётся корректный снимок векторных
часов. В тестовой конфигурации использовались:
\begin{itemize}[leftmargin=1.25cm]
    \item \texttt{10} реплик;
    \item \texttt{10} объектов;
    \item \texttt{100000} операций на объект.
\end{itemize}

Итоговый объём составил около \texttt{1\,000\,000} операций, а по данным
таблицы фактически было сохранено \texttt{\~1\,190\,000} строк. Это означает,
что в экспериментальном контуре таблица уже содержала объём порядка одного
миллиона записей, достаточный для первичных сравнительных замеров.

В классе \texttt{MainYt} реализованы два основных режима работы:
\begin{itemize}[leftmargin=1.25cm]
    \item генерация и пакетная запись операций;
    \item запуск набора простых benchmark-сценариев чтения.
\end{itemize}

Проверялись следующие запросы:
\begin{itemize}[leftmargin=1.25cm]
    \item поиск всех операций объекта по \texttt{objectId};
    \item поиск операций объекта от конкретной реплики;
    \item поиск операций объекта, реплики и транзакции;
    \item поиск операций объекта по префиксу пути \texttt{action.path} для поля \texttt{title};
    \item аналогичный поиск по пути \texttt{participants}.
\end{itemize}

Важно, что первые три сценария естественно согласованы с порядком ключа и
поэтому соответствуют модели хранения. Напротив, фильтрация по
\texttt{action.path} в текущей реализации не переносится в хранилище:
сначала считываются все операции объекта, а затем выполняется фильтрация в
памяти Java-приложения. Следовательно, эти сценарии показывают не столько
возможности индексного доступа YTsaurus, сколько стоимость полного чтения
истории объекта с последующей прикладной обработкой.

\withnumber{Измерения и результаты}

\mysubsection{Набор команд для анализа}

Для PostgreSQL и MongoDB в репозитории сохранены отдельные файлы с фактически использованными
командами анализа:
\begin{itemize}[leftmargin=1.25cm]
    \item \path{postgresql/resources/explain_analyze_commands.md};
    \item \path{mongodb/resources/explain_execution_stats_commands.md}.
\end{itemize}

Эти файлы фиксируют не только формулировки запросов, но и конкретные параметры,
использованные при локальных замерах.

\mysubsection{Результаты PostgreSQL}

Для PostgreSQL использовался датасет следующего объёма:
\begin{itemize}[leftmargin=1.25cm]
    \item \texttt{operations}: \texttt{1,000,000};
    \item \texttt{objects}: \texttt{8};
    \item \texttt{replicas}: \texttt{3}.
\end{itemize}

Дополнительно была снята статистика физического размера таблиц. Основная
нагрузка практически полностью сосредоточена в таблице \texttt{operations}:
её полный размер составил \texttt{393 MB}, из которых \texttt{221 MB}
приходится на данные таблицы и \texttt{171 MB} на индексы. Таблицы
\texttt{objects} и \texttt{replicas} при таком датасете занимают пренебрежимо
малый объём --- по \texttt{32 kB} каждая. Это хорошо согласуется с самой
моделью хранения: почти вся стоимость хранения в реляционном варианте связана
именно с большим журналом операций и поддерживающими его индексами.

Для измерений были выбраны следующие репрезентативные значения:
\begin{itemize}[leftmargin=1.25cm]
    \item \texttt{object\_id = 'event-7'}: \texttt{126,445} операций;
    \item \texttt{replica\_id = 'R2'}: \texttt{334,323} операций;
    \item \texttt{tx\_id = 'R3:tx:238257'};
    \item \texttt{op\_id = 'R3:238257'}.
\end{itemize}

Сводные результаты измерений для PostgreSQL вынесены в
\textbf{Приложение А}.

Показательный пример команды для чтения полной истории объекта:

\begin{lstlisting}[language=SQL,caption={Пример EXPLAIN ANALYZE для истории объекта в PostgreSQL}]
EXPLAIN (ANALYZE, BUFFERS)
SELECT
    op_id,
    tx_id,
    object_id,
    replica_id,
    clock,
    action
FROM operations
WHERE object_id = 'event-7'
ORDER BY op_id;
\end{lstlisting}

Результаты измерений показывают, что PostgreSQL особенно эффективен в точечных
и умеренно селективных сценариях. При этом широкие выборки и сортировка больших
подмножеств журнала остаются основным источником стоимости.

\mysubsection{Результаты MongoDB}

Для MongoDB использовался набор данных со следующими характеристиками:
\begin{itemize}[leftmargin=1.25cm]
    \item количество документов в \texttt{objects}: \texttt{10000};
    \item минимальное число операций в документе: \texttt{12};
    \item максимальное число операций в документе: \texttt{1000};
    \item среднее число операций в документе: \texttt{507.7832};
    \item оценка суммарного числа встроенных операций: примерно \texttt{5,077,832}.
\end{itemize}

По статистике \texttt{collStats} логический размер коллекции \texttt{objects}
составил \texttt{1165.44 MiB}, тогда как физический объём хранения
(\texttt{storageSize}) оказался равен \texttt{223.69 MiB}. Суммарный размер
индексов составил \texttt{57.85 MiB}, а общий объём коллекции вместе с
индексами --- \texttt{281.54 MiB}. Средний размер одного документа оказался
равен примерно \texttt{119.34 KiB}. Эти значения показывают, что встроенное
хранение длинных историй в документах приводит к значительному логическому
объёму данных, однако физически коллекция остаётся достаточно компактной за
счёт внутреннего сжатия движка хранения.

Для измерений использовались значения:
\begin{itemize}[leftmargin=1.25cm]
    \item \texttt{objectId = "event-3299"};
    \item \texttt{opId = "R1:365809"};
    \item \texttt{replicaId = "R5"};
    \item \texttt{action.type = "field.set"}.
\end{itemize}

Сводные результаты измерений для MongoDB вынесены в
\textbf{Приложение Б}.

Показательный пример команды для поиска объекта по вложенному \texttt{opId}:

\begin{lstlisting}[language=JavaScript,caption={Пример explain("executionStats") для MongoDB}]
db.objects.find(
  { "operations.opId": "R1:365809" },
  { _id: 0, objectId: 1 }
).explain("executionStats")
\end{lstlisting}

Полученные результаты подтверждают, что MongoDB лучше всего проявляет себя в
объектно-центричных сценариях: поиск объекта, чтение полной истории и доступ к
последним операциям выполняются естественно и быстро, тогда как глобальные
выборки по операциям остаются менее характерными для этой модели.

\mysubsection{Результаты YTsaurus}

Согласно сохранённому логу выполнения, были получены следующие времена:
\begin{itemize}[leftmargin=1.25cm]
    \item \texttt{findByObjectId}: \texttt{5357 ms}, \texttt{100000} строк;
    \item \texttt{findByObjectAndReplica}: \texttt{714 ms}, \texttt{9957} строк;
    \item \texttt{findByObjectReplicaAndTx}: \texttt{136 ms}, \texttt{1} строка;
    \item \texttt{findByObjectIdAndPath(title)}: \texttt{18114 ms}, \texttt{8524} строки;
    \item \texttt{findByObjectIdAndPath(participants)}: \texttt{19152 ms}, \texttt{30237} строк.
\end{itemize}

Наблюдаемая картина логична. Чем точнее запрос совпадает с префиксом
составного ключа, тем быстрее он выполняется. Самым дешёвым оказывается
сценарий чтения одной операции по сочетанию \texttt{objectId},
\texttt{replicaId} и \texttt{txId}. Напротив, сценарии по \texttt{path}
получаются существенно дороже, поскольку требуют сначала извлечь всю историю
объекта, а затем отфильтровать её на прикладном уровне. По сути, это
подтверждает, что без отдельного проектирования ключа или дополнительного
представления таблицы запросы по вложенному пути действия для YTsaurus в такой
схеме оказываются неестественными.

\mysubsection{Параметры таблицы и объём хранения}

Скриншоты интерфейса YTsaurus позволяют зафиксировать важные параметры
физического хранения таблицы \path{//home/dev/n.d.ivanov/operations}.

\illustration{../YTsaurus/yt_1.jpg}{0.92}{Общие параметры таблицы YTsaurus}

Из общих атрибутов видно, что таблица размещена в аккаунте \texttt{dev},
использует \texttt{primary medium = default}, \texttt{tablet cell bundle = default},
а режим \texttt{in-memory} отключён. Это означает, что в данном эксперименте
таблица не была специально оптимизирована под хранение данных в памяти и
оценивалась в более консервативной конфигурации.

\illustration{../YTsaurus/YT_2.jpg}{0.92}{Параметры хранения таблицы YTsaurus}

На уровне storage-настроек таблица была сконфигурирована как
\texttt{optimize for lookup}, с кодеком сжатия \texttt{lz4} и репликацией
\texttt{3}. Такое сочетание выглядит оправданным для сценария point lookup и
чтения по префиксу ключа: выбран режим, ориентированный именно на запросы по
ключу, а не на последовательный scan большого массива данных.

\illustration{../YTsaurus/YT_3.jpg}{0.55}{Статистика объёма таблицы YTsaurus}

По данным интерфейса таблица содержала примерно \texttt{1\,190\,000} строк,
\texttt{3} chunk'а, \texttt{457.44 MiB} в несжатом виде и \texttt{111.68 MiB}
в сжатом виде. При трёхкратной репликации суммарное занимаемое дисковое
пространство составило \texttt{336.20 MiB}, а \texttt{data weight} был равен
\texttt{365.63 MiB}.

Эти значения позволяют сделать несколько выводов. Во-первых, сжатие
\texttt{lz4} оказалось весьма эффективным: отношение несжатого объёма к
сжатому составляет примерно \texttt{4.1x}. Во-вторых, итоговый показатель
\texttt{total disk space} практически совпадает с тройным объёмом сжатых
данных, что хорошо согласуется с включённой репликацией \texttt{3}. В-третьих,
даже при хранении около 1.19 млн операций суммарный физический объём таблицы
остаётся умеренным, что делает подобную схему хранения реалистичной для
журнала CRDT-операций, по крайней мере на масштабе миллионов записей.

\mysubsection{Промежуточный вывод}

Проведённый эксперимент показывает, что YTsaurus может использоваться как
распределённое хранилище append-only журнала \texttt{op-based CRDT}-операций,
если ключ проектируется под наиболее частые сценарии чтения. В такой схеме
платформа естественно поддерживает выборки по \texttt{objectId},
\texttt{replicaId} и \texttt{txId}, однако запросы по содержимому вложенного
\texttt{action.path} без дополнительной перестройки модели оказываются заметно
менее эффективными. При этом отсутствие доступа к \texttt{EXPLAIN}-подобным
механизмам в используемом кластере ограничивает анализ уровнем внешнего
времени исполнения и не позволяет детально разобрать внутреннюю стратегию
выполнения запросов.

\withoutnumber{Заключение}

В работе были рассмотрены три варианта хранения \texttt{op-based CRDT}-операций:
реляционный на PostgreSQL, документный на MongoDB и распределённый табличный
вариант на YTsaurus. Для каждого из подходов была проанализирована единица
хранения, структура данных, характер типовых запросов и практические ограничения,
возникающие при работе с журналом операций.

PostgreSQL показал себя как наиболее естественный выбор для модели глобального
immutable operation log, в которой операция рассматривается как самостоятельная
запись, а важную роль играют ссылочная целостность, прозрачность схемы и
возможность детального анализа выполнения запросов. MongoDB, напротив, оказался
более естественным для объектно-центричного хранения, где история операций
агрегируется внутри документа и может быть быстро прочитана целиком для одного
объекта. YTsaurus в исследованной постановке продемонстрировал, что также может
использоваться для хранения большого append-only журнала операций, если
составной ключ проектируется в соответствии с реальными сценариями доступа.

Проведённое сравнение показывает, что выбор хранилища для \texttt{op-based CRDT}
не сводится только к сравнению производительности отдельных запросов. На
практике определяющими оказываются архитектурная модель хранения, характер
доступа к истории операций, требования к масштабируемости и доступные средства
наблюдаемости. В этом смысле PostgreSQL, MongoDB и YTsaurus представляют три
разных подхода к хранению одной и той же прикладной структуры данных, и каждый
из них оказывается оправданным в собственном классе сценариев.

Если рассматривать единицу хранения, то PostgreSQL и YTsaurus в данной работе
оказались ближе друг к другу, чем к MongoDB. В обоих случаях операция является
самостоятельной записью, пригодной для выборок по идентификаторам объекта,
реплики и транзакции. Однако различие между ними принципиально состоит в том,
что PostgreSQL даёт развитые механизмы реляционного анализа, прозрачные планы
выполнения запросов и привычную схему ограничений целостности, тогда как
YTsaurus сильнее ориентирован на распределённое хранение и доступ по ключу в
масштабируемой инфраструктурной среде.

MongoDB, напротив, реализует иную модель: операция не существует как отдельный
top-level элемент хранения, а включается во встроенный массив внутри документа
объекта. Это даёт естественное преимущество в сценариях чтения полной истории
одного объекта, но делает менее удобными глобальные выборки по операциям через
весь набор данных.

С точки зрения запросов различие также проявилось достаточно явно. PostgreSQL
демонстрирует сильные результаты в точечных и умеренно селективных сценариях,
а также предоставляет детальную диагностику через \texttt{EXPLAIN ANALYZE}.
MongoDB хорошо поддерживает доступ к истории конкретного объекта и поиск по
вложенным полям операций за счёт multikey-индексов. YTsaurus в исследованной
конфигурации показал, что лучше всего подходит для запросов, согласованных с
порядком составного ключа, то есть прежде всего по \texttt{objectId},
\texttt{replicaId} и \texttt{txId}, тогда как фильтрация по содержимому
\texttt{action.path} без специального проектирования модели оказывается
существенно менее эффективной.

Отдельно следует отметить различие в глубине экспериментального анализа.
Для PostgreSQL и MongoDB удалось исследовать не только время, но и сами планы
выполнения запросов. Для YTsaurus на доступном кластере VK такой возможности
не было: измерялись только внешние времена исполнения и размеры результатов.
Поэтому выводы по YTsaurus неизбежно имеют более прикладной и менее
плано-ориентированный характер.

При этом с точки зрения самой прикладной логики работы с композитными
\texttt{CRDT}-объектами наиболее естественным в рамках данного исследования
оказался именно MongoDB-подход. Причина состоит в том, что в документной модели
исследователь и разработчик оперируют прежде всего объектом как целостной
 сущностью, а не множеством отдельных строк журнала. История изменений,
 относящаяся к одному \texttt{objectId}, оказывается собрана в пределах одного
 документа и потому воспринимается как единый агрегат, близкий к исходному
 прикладному представлению календарного события.

Такой способ хранения лучше соответствует интуитивной модели предметной области.
При работе с operation-based \texttt{CRDT} на прикладном уровне обычно важно
мыслить не отдельными операциями как таковыми, а состоянием и эволюцией
конкретного объекта, над которым эти операции выполняются. В MongoDB эта
перспектива поддерживается самой структурой хранилища: объект и его история
оказываются связаны непосредственно, без необходимости каждый раз восстанавливать
эту связь через внешний журнал операций. В результате модель оказывается проще
для концептуального восприятия, ближе к объектному устройству приложения и
удобнее для сценариев, где основной единицей анализа является именно отдельный
CRDT-объект.

Именно поэтому, хотя PostgreSQL предоставляет более сильные средства анализа,
строгую схему и удобную работу с глобальным operation log, а YTsaurus интересен
как распределённая масштабируемая платформа, MongoDB в данной постановке можно
считать наиболее естественным вариантом с точки зрения представления самой
предметной сущности. Он лучше отражает объектную природу исследуемых данных и
позволяет рассматривать историю операций не как внешнее множество событий, а
как внутреннюю характеристику конкретного объекта.

\withoutnumber{Приложение А. Сводные результаты PostgreSQL}

\begin{landscape}
\begin{reporttable}{
\tablecaption{Ключевые результаты EXPLAIN ANALYZE для PostgreSQL}{postgresql-measurements}
\begin{tabular}{|p{5.0cm}|p{8.0cm}|p{3.0cm}|p{8.5cm}|}
\hline
\textbf{Сценарий} & \textbf{План} & \textbf{Время} & \textbf{Наблюдение} \\
\hline
Размер таблицы \texttt{operations} & Таблица \texttt{operations} + индексы & \texttt{393 MB} & \texttt{221 MB} данных и \texttt{171 MB} индексов; основной объём хранения сосредоточен в журнале операций \\
\hline
История объекта по \texttt{object\_id} & \texttt{Bitmap Index Scan} + \texttt{Parallel Bitmap Heap Scan} + \texttt{Gather Merge} & \texttt{\~573.7 ms} & Индекс по \texttt{object\_id} используется, но значимая часть стоимости уходит в чтение heap pages и сортировку \\
\hline
Операции транзакции по \texttt{tx\_id} & \texttt{Index Scan} по \texttt{idx\_operations\_tx\_id} & \texttt{\~0.245 ms} & Быстрый точечный сценарий с эффективным использованием B-tree индекса \\
\hline
Операции реплики по \texttt{replica\_id} & \texttt{Parallel Seq Scan} + сортировка & \texttt{\~210.0 ms} & Индекс не выбран из-за низкой селективности условия \\
\hline
Поиск по \texttt{op\_id} & \texttt{Index Scan} по \texttt{operations\_pkey} & \texttt{\~0.213 ms} & Первичный ключ полностью закрывает данный сценарий \\
\hline
Фильтрация по частому \texttt{action.type = field.set} & \texttt{Parallel Seq Scan} & \texttt{\~1004.4 ms} & Индекс невыгоден для очень частого значения \\
\hline
Фильтрация по селективному \texttt{action.type = set.remove} & \texttt{Bitmap Index Scan} + \texttt{Parallel Bitmap Heap Scan} & \texttt{\~275.4 ms} & Индекс по типу действия работает при селективном значении \\
\hline
Поиск по фрагменту \texttt{action} & \texttt{Bitmap Index Scan} по GIN & \texttt{\~485.2 ms} & GIN-индекс полезен для containment-запросов по \texttt{JSONB} \\
\hline
Поиск по фрагменту \texttt{clock} & \texttt{Bitmap Index Scan} по GIN & \texttt{\~0.534 ms} & Индекс по \texttt{clock} реально используется и эффективен \\
\hline
\end{tabular}
}
\end{reporttable}
\end{landscape}

\withoutnumber{Приложение Б. Сводные результаты MongoDB}

\begin{landscape}
\begin{reporttable}{
\tablecaption{Ключевые результаты executionStats для MongoDB}{mongodb-measurements}
\begin{tabular}{|p{5.2cm}|p{7.0cm}|p{4.0cm}|p{8.0cm}|}
\hline
\textbf{Сценарий} & \textbf{План} & \textbf{Показатели} & \textbf{Наблюдение} \\
\hline
Размер коллекции \texttt{objects} & \texttt{collStats} & \texttt{1165.44 MiB} логически, \texttt{223.69 MiB} storage, \texttt{57.85 MiB} индексы & При \texttt{10000} документах и среднем размере документа \texttt{119.34 KiB} физический объём остаётся умеренным за счёт сжатия \\
\hline
Поиск объекта по \texttt{objectId} & \texttt{IXSCAN -> FETCH} & \texttt{1 key / 1 doc} & Базовый сценарий полностью поддержан уникальным индексом \\
\hline
Чтение полной истории объекта & \texttt{IXSCAN -> FETCH} & \texttt{1 key / 1 doc}, \texttt{2 ms} & История объекта читается как один документ \\
\hline
Поиск по \texttt{operations.opId} & \texttt{IXSCAN -> FETCH -> PROJECTION} & \texttt{1 key / 1 doc} & Multikey индекс работает эффективно \\
\hline
Поиск по \texttt{operations.replicaId} & \texttt{IXSCAN -> FETCH -> LIMIT} & \texttt{100 keys / 100 docs} & Индекс используется корректно, но сценарий вторичен для embedded-модели \\
\hline
Поиск по \texttt{operations.action.type} & \texttt{IXSCAN -> FETCH -> LIMIT} & \texttt{100 keys / 100 docs} & Индекс по типу действия работает корректно \\
\hline
Чтение последних операций через \texttt{\$slice} & \texttt{IXSCAN -> FETCH -> PROJECTION} & \texttt{1 key / 1 doc}, \texttt{4 ms} & Сценарий хорошо соответствует объектной модели \\
\hline
\end{tabular}
}
\end{reporttable}
\end{landscape}

\withoutnumber{Приложение В. Сводные результаты YTsaurus}

\begin{landscape}
\begin{reporttable}{
\tablecaption{Ключевые результаты измерений для YTsaurus}{ytsaurus-measurements}
\begin{tabular}{|p{5.3cm}|p{5.0cm}|p{3.0cm}|p{3.0cm}|p{9.0cm}|}
\hline
\textbf{Сценарий} & \textbf{Условие} & \textbf{Время} & \textbf{Результат} & \textbf{Наблюдение} \\
\hline
Все операции объекта & \texttt{objectId = event-0} & \texttt{5357 ms} & \texttt{100000} строк & Чтение большой истории одного объекта по первому столбцу составного ключа; сценарий естественен для выбранной схемы, но остаётся дорогим из-за объёма результата \\
\hline
Операции объекта конкретной реплики & \texttt{objectId = event-0}, \texttt{replicaId = replica-1} & \texttt{714 ms} & \texttt{9957} строк & Существенно быстрее за счёт более точного совпадения с префиксом ключа \\
\hline
Операция по объекту, реплике и транзакции & \texttt{objectId = event-0}, \texttt{replicaId = replica-1}, \texttt{txId = replica-1:tx:1000} & \texttt{136 ms} & \texttt{1} строка & Наиболее селективный сценарий; лучше всего соответствует key-lookup модели таблицы \\
\hline
Операции объекта по пути \texttt{title} & \texttt{objectId = event-0}, фильтрация по \texttt{action.path} & \texttt{18114 ms} & \texttt{8524} строки & Поиск по пути не поддерживается напрямую ключом: сначала считывается вся история объекта, затем выполняется фильтрация в приложении \\
\hline
Операции объекта по пути \texttt{participants} & \texttt{objectId = event-0}, фильтрация по \texttt{action.path} & \texttt{19152 ms} & \texttt{30237} строк & Самый дорогой сценарий из-за полного чтения истории и последующей прикладной фильтрации по пути \\
\hline
\end{tabular}
}
\end{reporttable}
\end{landscape}

\withoutnumber{Приложение Г. Сравнение подходов хранения}

\begin{landscape}
\begin{reporttable}{
\tablecaption{Сравнение вариантов хранения}{storage-comparison}
\begin{tabular}{|p{4.4cm}|p{7.0cm}|p{7.0cm}|p{6.0cm}|}
\hline
\textbf{Критерий} & \textbf{PostgreSQL} & \textbf{MongoDB} & \textbf{YTsaurus} \\
\hline
Единица хранения & Отдельная операция & Документ объекта со встроенной историей & Отдельная операция в distributed dynamic table \\
\hline
Объём хранения & \texttt{operations}: \texttt{393 MB} total, из них \texttt{221 MB} данные и \texttt{171 MB} индексы & \texttt{objects}: \texttt{1165.44 MiB} logical size, \texttt{223.69 MiB} storage, \texttt{57.85 MiB} индексы & \texttt{\~1\,190\,000} строк, \texttt{457.44 MiB} uncompressed, \texttt{111.68 MiB} compressed, \texttt{336.20 MiB} total disk space \\
\hline
Физическая структура & Таблицы \texttt{replicas}, \texttt{objects}, \texttt{operations} & Коллекция \texttt{objects} и массив \texttt{operations} & Dynamic table \path{//home/dev/n.d.ivanov/operations} с составным ключом \texttt{objectId, replicaId, txId, opId} \\
\hline
Сильная сторона & Глобальный журнал операций, аналитика, ссылочная целостность & Быстрое чтение полной истории объекта & Масштабируемое распределённое хранение и эффективный lookup по префиксу ключа \\
\hline
Слабая сторона & Дорогая сортировка больших выборок, часть структуры в \texttt{JSONB} & Ограничение размера документа и менее естественные глобальные выборки & Ограниченная наблюдаемость плана выполнения и слабая естественность запросов по \texttt{action.path} без отдельного проектирования модели \\
\hline
Инструмент анализа & \texttt{EXPLAIN ANALYZE} & \texttt{explain("executionStats")} & Замер внешнего времени выполнения без \texttt{EXPLAIN}-API в доступном контуре \\
\hline
\end{tabular}
}
\end{reporttable}
\end{landscape}

\end{document}
